% GENERATED FILE. Edit canonical lessons, then run npm run generate:books.
% canonical-source-hash: fnv1a64:c90b1b4a3f39da63
% canonical-lessons: MR-C81-thambne, MR-C81-shodhne, MR-C81-haravne, MR-C81-ghalne, MR-C81-sodne

\chapter{To stop, To look for, To lose, To put on, To leave}
\label{ch:mr-to-stop-to-look-for-to-lose-to-put-on-to-leave}
\hlchaptermodality{81}

\begin{chapteropening}
\textbf{By the end of this chapter:} \emph{I can say to stop, to wait, to look for, to lose, to put on and to leave, to let go.}

Complete the last lesson of chapter 81: I can say to stop, to wait, to look for, to lose, to put on and to leave, to let go.
\end{chapteropening}

\section[\texorpdfstring{\mr{थांबणे} (thāmbṇe)}{थांबणे (thāmbṇe)}]{\mr{थांबणे} (thāmbṇe) — to stop, to wait}

\label{lesson:MR-C81-thambne}



\begin{quote}
Before the new one: say the Marathi for to catch, then the Marathi for to cut.
\end{quote}

\subsection*{You'll want to know: \mr{थांबणे}}
\textbf{\mr{थांबणे}} — \emph{thāmbṇe} — \textquotedblleft{}to stop, to wait\textquotedblright{}.

\begin{cousinweb}
\textbf{\mr{थांबणे}} (\emph{thāmbṇe}) is a tadbhava, a word inherited from Sanskrit and reshaped by sound change. It comes from \emph{stambhate}, \textquotedblleft{}supports, arrests\textquotedblright{} (CDIAL 13683). CDIAL glosses the Marathi verb \textquotedblleft{}to stop\textquotedblright{}.
\end{cousinweb}

\begin{grammarlens}[title={What you've built}]
One more word for this chapter.
\end{grammarlens}

\subsection*{Your turn}
\begin{itemize}
\raggedright
  \item \emph{Say it:} \emph{thāmbṇe}
  \item \emph{Say it:} \emph{thāmbṇe}, once more
  \item \emph{Say it:} say \emph{thāmbṇe}
  \item \emph{Recall:} say the Marathi for to catch, then the Marathi for to cut, then say \emph{thāmbṇe} again
\end{itemize}

\begin{tcolorbox}[breakable,colback=teal!4,colframe=teal!35!black,title={Before you move on}]
What does \mr{थांबणे} mean? (\textquotedblleft{}To stop, to wait\textquotedblright{}.) Say it once more.
\end{tcolorbox}
\section[\texorpdfstring{\mr{शोधणे} (shodhṇe)}{शोधणे (shodhṇe)}]{\mr{शोधणे} (shodhṇe) — to look for}

\label{lesson:MR-C81-shodhne}



\begin{quote}
Before the new one: say the Marathi for to cut, then the Marathi for to stop.
\end{quote}

\subsection*{You'll want to know: \mr{शोधणे}}
\textbf{\mr{शोधणे}} — \emph{shodhṇe} — \textquotedblleft{}to look for\textquotedblright{}.

\begin{cousinweb}
\textbf{\mr{शोधणे}} (\emph{shodhṇe}) is a tadbhava, a word inherited from Sanskrit and reshaped by sound change. It comes from \emph{śoddhum}, \textquotedblleft{}to purify\textquotedblright{} (CDIAL 12626). CDIAL glosses the Marathi verb \textquotedblleft{}to purify, search\textquotedblright{}.
\end{cousinweb}

\begin{grammarlens}[title={What you've built}]
One more word for this chapter.
\end{grammarlens}

\subsection*{Your turn}
\begin{itemize}
\raggedright
  \item \emph{Say it:} \emph{shodhṇe}
  \item \emph{Say it:} \emph{shodhṇe}, once more
  \item \emph{Say it:} say \emph{shodhṇe}
  \item \emph{Recall:} say the Marathi for to cut, then the Marathi for to stop, then say \emph{shodhṇe} again
\end{itemize}

\begin{tcolorbox}[breakable,colback=teal!4,colframe=teal!35!black,title={Before you move on}]
What does \mr{शोधणे} mean? (\textquotedblleft{}To look for\textquotedblright{}.) Say it once more.
\end{tcolorbox}
\section[\texorpdfstring{\mr{हरवणे} (haravṇe)}{हरवणे (haravṇe)}]{\mr{हरवणे} (haravṇe) — to lose}

\label{lesson:MR-C81-haravne}



\begin{quote}
Before the new one: say the Marathi for to stop, then the Marathi for to look for.
\end{quote}

\subsection*{You'll want to know: \mr{हरवणे}}
\textbf{\mr{हरवणे}} — \emph{haravṇe} — \textquotedblleft{}to lose\textquotedblright{}.

\begin{cousinweb}
\textbf{\mr{हरवणे}} (\emph{haravṇe}) is a tadbhava, a word inherited from Sanskrit and reshaped by sound change. It comes from \emph{hārayati}, \textquotedblleft{}causes to be taken; loses\textquotedblright{} (CDIAL 14061). CDIAL lists the Marathi form \emph{haraviṇẽ}, \textquotedblleft{}to lose, cast away\textquotedblright{}.
\end{cousinweb}

\begin{grammarlens}[title={What you've built}]
One more word for this chapter.
\end{grammarlens}

\subsection*{Your turn}
\begin{itemize}
\raggedright
  \item \emph{Say it:} \emph{haravṇe}
  \item \emph{Say it:} \emph{haravṇe}, once more
  \item \emph{Say it:} say \emph{haravṇe}
  \item \emph{Recall:} say the Marathi for to stop, then the Marathi for to look for, then say \emph{haravṇe} again
\end{itemize}

\begin{tcolorbox}[breakable,colback=teal!4,colframe=teal!35!black,title={Before you move on}]
What does \mr{हरवणे} mean? (\textquotedblleft{}To lose\textquotedblright{}.) Say it once more.
\end{tcolorbox}
\section[\texorpdfstring{\mr{घालणे} (ghālṇe)}{घालणे (ghālṇe)}]{\mr{घालणे} (ghālṇe) — to put on (clothes)}

\label{lesson:MR-C81-ghalne}



\begin{quote}
Before the new one: say the Marathi for to look for, then the Marathi for to lose.
\end{quote}

\subsection*{You'll want to know: \mr{घालणे}}
\textbf{\mr{घालणे}} — \emph{ghālṇe} — \textquotedblleft{}to put on (clothes)\textquotedblright{}.

\begin{cousinweb}
\textbf{\mr{घालणे}} (\emph{ghālṇe}) is a tadbhava, a word inherited and reshaped by sound change. It comes from \emph{*ghalyati}, \textquotedblleft{}pours\textquotedblright{}, a form CDIAL reconstructs from its descendants (CDIAL 4457). CDIAL glosses the Marathi verb \textquotedblleft{}to pour\textquotedblright{}.
\end{cousinweb}

\begin{grammarlens}[title={What you've built}]
One more word for this chapter.
\end{grammarlens}

\subsection*{Your turn}
\begin{itemize}
\raggedright
  \item \emph{Say it:} \emph{ghālṇe}
  \item \emph{Say it:} \emph{ghālṇe}, once more
  \item \emph{Say it:} say \emph{ghālṇe}
  \item \emph{Recall:} say the Marathi for to look for, then the Marathi for to lose, then say \emph{ghālṇe} again
\end{itemize}

\begin{tcolorbox}[breakable,colback=teal!4,colframe=teal!35!black,title={Before you move on}]
What does \mr{घालणे} mean? (\textquotedblleft{}To put on (clothes)\textquotedblright{}.) Say it once more.
\end{tcolorbox}
\section[\texorpdfstring{\mr{सोडणे} (soḍṇe)}{सोडणे (soḍṇe)}]{\mr{सोडणे} (soḍṇe) — to leave, to let go}

\label{lesson:MR-C81-sodne}



\begin{quote}
Before the new one: say the Marathi for to lose, then the Marathi for to put on.
\end{quote}

\subsection*{You'll want to know: \mr{सोडणे}}
\textbf{\mr{सोडणे}} — \emph{soḍṇe} — \textquotedblleft{}to leave, to let go\textquotedblright{}.

\begin{cousinweb}
\textbf{\mr{सोडणे}} (\emph{soḍṇe}) is a tadbhava, a word inherited from Sanskrit and reshaped by sound change. It comes from \emph{kṣoṭayati}, \textquotedblleft{}throws\textquotedblright{} (CDIAL 3747).
\end{cousinweb}

\begin{grammarlens}[title={What you've built}]
One more word for this chapter.
\end{grammarlens}

\subsection*{Your turn}
\begin{itemize}
\raggedright
  \item \emph{Say it:} \emph{soḍṇe}
  \item \emph{Say it:} \emph{soḍṇe}, once more
  \item \emph{Say it:} say \emph{soḍṇe}
  \item \emph{Recall:} say the Marathi for to lose, then the Marathi for to put on, then say \emph{soḍṇe} again
\end{itemize}

\begin{tcolorbox}[breakable,colback=teal!4,colframe=teal!35!black,title={Before you move on}]
What does \mr{सोडणे} mean? (\textquotedblleft{}To leave, to let go\textquotedblright{}.) Say it once more.
\end{tcolorbox}
